%% ============================================================
%% Point-by-point response to the Academic Editor
%% LegalShadow · Engineering Proceedings (MDPI) · ICI2ST 2026 volume
%% Standalone: compile with pdflatex. No external class required.
%% AUTOR: verificar los números de sección/página contra el PDF final
%%        antes de enviar (se indican con \S y con "p." aproximada).
%% ============================================================
\documentclass[11pt,a4paper]{article}

\usepackage[T1]{fontenc}
\usepackage[utf8]{inputenc}
\usepackage{lmodern}
\usepackage[english]{babel}
\usepackage[left=2.3cm,right=2.3cm,top=2.3cm,bottom=2.3cm]{geometry}
\usepackage{microtype}
\usepackage{enumitem}
\usepackage{xcolor}
\usepackage{amsmath}
\usepackage{mdframed}
\usepackage[hidelinks]{hyperref}

\definecolor{cmtbg}{HTML}{F1F5F9}
\definecolor{cmtrule}{HTML}{94A3B8}

\newenvironment{comment}[1]
  {\par\vspace{6pt}\noindent\textbf{#1}\par\vspace{2pt}%
   \begin{mdframed}[backgroundcolor=cmtbg, linecolor=cmtrule, linewidth=0.6pt,
     innertopmargin=5pt, innerbottommargin=5pt, innerleftmargin=7pt,
     innerrightmargin=7pt, skipabove=2pt, skipbelow=5pt]\itshape}
  {\end{mdframed}}

\newcommand{\respuesta}{\par\noindent\textbf{Response.}\ }
\newcommand{\changed}[1]{\par\vspace{3pt}\noindent\textit{Changes in the manuscript:} #1\par}
\setlist[itemize]{leftmargin=1.3em, itemsep=2pt, topsep=3pt}
\pagestyle{plain}

\begin{document}

\begin{center}
{\Large\bfseries Response to the Academic Editor}\\[4pt]
{\large \emph{LegalShadow: An Ontology-Guided Digital Shadow Architecture for\\ Semantic Validation of Judicial Digital Ecosystems}}\\[6pt]
Patricio M. Paccha-Angamarca, Erwin J. Sacoto-Cabrera, V\'ictor V. Velepucha-Bonett\\[4pt]
\emph{Engineering Proceedings} (MDPI) --- ICI2ST 2026 volume --- \today
\end{center}

\vspace{6pt}

\noindent Dear Academic Editor,

\vspace{4pt}

\noindent Thank you for an unusually constructive report. The five major points identified real weaknesses, and addressing them has improved the paper in ways we had not anticipated: the delimitation you asked for in \S3.3 forced us to state the direction of validation precisely, and that turned out to be the sharpest way to express what the contribution is. We have revised the manuscript accordingly and respond to each point below. All changes are marked in the source with \texttt{\% REV2} comments; the revised manuscript runs to 14 pages including references, against 10 for the previous version, the growth being almost entirely the material you requested.

\vspace{4pt}

\noindent A summary of what changed:

\begin{itemize}
  \item A new Section~2, \emph{Related Work and Delimitation}, with a comparison table and nine new third-party references (\S3.3).
  \item A new subsection documenting the anchoring $\Gamma$ as an auditable rule table, plus a JSON-LD listing of a real shadow component (\S3.2, \S3.4).
  \item A fully worked $\Psi$ calculation, from HTTP evidence to remediation priority (\S3.4).
  \item An expanded reflexive finding, split into a methodological consequence and a concrete publication agenda (\S3.5).
  \item A new discussion block on external validity with a four-step transfer protocol (\S3.1).
  \item Minor: over-claims softened, terminology fixed, all figures and tables now described in the text (\S4).
\end{itemize}

\hrulefill

\section*{3.1 Limited empirical scope and external validity}

\begin{comment}{Editor's comment}
The study is restricted to a single institution (Ecuadorian Judiciary Council) over a relatively short observation window (17 runs). The authors must discuss more carefully the limits of generalisation and outline how the method could be transferred to other judicial ecosystems.
\end{comment}

\respuesta We accept this fully. The previous version compressed the issue into a single clause of the Limitations paragraph, which was not proportionate. We have added a dedicated block, \emph{External validity and how to transfer the method}, at the end of the Discussion.

Our response separates three claims that the previous text allowed a reader to conflate. The \emph{numerical} result generalises to nothing: 46/100 and $\Psi_{\mathrm{global}}=34.9\%$ describe one institution in August 2026. The \emph{diagnostic pattern}---high declarative maturity alongside low semantic coverage concentrated in open data---is stated explicitly as a hypothesis that a single case makes plausible and cannot test; we note that the open government data literature reporting a publication-versus-reuse gap across hundreds of portals makes it unlikely to be idiosyncratic, but we label that inference by analogy rather than evidence. Only the \emph{instrument} generalises.

On transfer, we now give a four-step protocol and, more usefully, name the step that actually costs something. Steps 1, 3 and 4 are cheap: the target URL is already a parameter, the ontology is international in eight of ten dimensions, and cross-institutional comparison uses $\Psi$ and $\delta$ rather than raw signal counts. Step 2 is the obstacle: roughly a fifth of the rules in $\Gamma$ key on national artefacts---the Ecuadorian open data catalogue, national interoperability schemes---and must be re-expressed against equivalent infrastructure elsewhere. We prefer to state this plainly rather than claim frictionless portability. We also now close the Conclusions by locating country-specificity precisely in the rule table rather than asserting that ``nothing in the architecture is country-specific''.

\changed{New paragraph pair \emph{External validity and how to transfer the method} in Section~6; Conclusions rewritten on this point; abstract now mentions the transfer protocol.}

\section*{3.2 Transparency of ontology scoring and component-to-class mapping}

\begin{comment}{Editor's comment}
The maturity scores that feed the model remain expert-elicited. The mapping $\Gamma$ (observed components $\rightarrow$ ontology classes) is critical for reproducibility, yet it is insufficiently detailed. More information is needed on how these mappings are decided and whether they can be independently verified.
\end{comment}

\respuesta This was the most consequential omission and we have addressed it in two places.

On $\Gamma$: a new subsection, \emph{The anchoring $\Gamma$ and how to audit it}, states what the mapping actually is. We also clarified something the previous version left implicit and which we now realise was the source of the opacity you identified: $\Gamma$ is not applied after the fact to whatever a crawl happened to collect. The ontology \emph{drives} the acquisition---each class declares the observable evidence that would materialise it, so the OWL file is the search plan---and $\Gamma$ is that same rule table read in the opposite direction. Section~4 and Algorithm~1 have been rewritten accordingly, with a new first step that derives the probe set from $\Omega$. The rules themselves are not inferred, learned or negotiated but a versioned table of the form ``if deterministic signal $s$ appears in component $v$, then $\mathrm{class}(s)\in\Gamma(v)$''. Every rule fires on a syntactic condition re-evaluable on the stored snapshot---an HTTP header, a DOM element, a URL pattern, a well-known path---and never on a human reading of a page. We give concrete instances (\texttt{OpenData\_Layer} requires a resolvable catalogue reference; \texttt{Transport\_Security} requires an HSTS header on the stored response) and state the three properties that follow: totality and referential closure by construction, separate auditability of observation and interpretation, and conservatism, since absence of signal is read as absence in the public surface and never as absence in the institution. The new JSON-LD listing (Figure~2) shows a real component with its evidence block, its \texttt{maltg:mapsTo} array and its recorded gap, so the reader can see the mapping rather than take our word for it.

On $\sigma$: rather than defend the elicitation, we now separate the two quantities by epistemic status. $\mathrm{onto}(d)$ is expert-elicited and carries judgement; $\Psi(d)$ is a deterministic function of the rule table and the snapshots and carries none. This is what makes \emph{partial} replication possible: a reviewer who rejects our rubric can still recompute every $\Psi(d)$ from the published artefacts, substitute their own $\mathrm{onto}$ vector and derive their own agenda. We add that the ordering is robust---the seven dimensions with $\Psi<0.35$ hold the top seven positions by $\delta$ under any $\mathrm{onto}$ assignment keeping AI, ITIL and security above 70---because $\delta$ is monotone in $\mathrm{onto}$ at fixed $\Psi$. The absence of an independent panel remains in Limitations, and independent elicitation of $\sigma$ is now named explicitly as future work.

\changed{New subsection on $\Gamma$ in Section~4; new Figure~2 (JSON-LD listing); new paragraph on the epistemic separation of $\mathrm{onto}$ and $\Psi$ in the Scoring protocol of Section~5.}

\section*{3.3 Novelty claim}

\begin{comment}{Editor's comment}
The assertion that no prior architecture simultaneously performs automatic shadow construction from public web evidence, JSON-LD serialisation anchored to a governance ontology, and hierarchical coverage validation needs stronger substantiation against related work in digital shadows for public administration, open-data maturity assessment, and semantic monitoring of government portals.
\end{comment}

\respuesta Agreed: the claim was asserted rather than argued. We have added Section~2, \emph{Related Work and Delimitation}, which reviews exactly the three lines you name and nine new third-party references, and closes with a comparison table along the three axes of the claim.

The substantiation is now positional rather than rhetorical. Digital twins of public systems are increasingly \emph{assessed}---Haraguchi et al.'s CITYSTEPS maturity model, Diaz-Sarachaga's 33-attribute Delphi rating for urban twins, Anshari et al.\ on twins as an e-government platform---but every one of these assesses a twin programme run by the system's own operator, using expert rubrics over operator-supplied information. Automated open government data assessment shares our acquisition philosophy---Neumaier et al.\ over 260 portals, Kubler et al.\ over 250 with AHP, Vetr\`o et al.\ at dataset level---but takes the dataset and its metadata record as unit of analysis and DCAT or FAIR as yardstick; the institution's architecture is out of scope. Ontology and knowledge-graph quality metrics resemble $\Psi$ formally---Zaitoun et al., Abad-Navarro and Mart\'inez-Costa, Seo et al.---but measure how well an \emph{ontology} covers a domain, or how well a graph is internally formed.

Stating it that way isolated what is actually new, and we are grateful for the push. Each ingredient exists separately; what we did not find is the combination, and specifically the \emph{direction} of validation: the ontology as fixed yardstick and the observed artefact as the thing scored. That is precisely the verification gap Karabulut et al.\ report as unaddressed. We have also removed the bare ``to the best of our knowledge'' and bounded the claim to the scope of the Section~2 comparison, as suggested in your minor comments.

\changed{New Section~2 with Table~1; novelty paragraph in Section~1 rewritten to point at it and to drop the unbounded phrasing; nine references added.}

\section*{3.4 Presentation and completeness of results}

\begin{comment}{Editor's comment}
Concrete illustrations of the JSON-LD shadow fragments and additional worked examples of the $\Psi$ calculation would significantly improve readability. The contrast between the four-dimension declarative maturity index and the ten-dimension semantic coverage should be made even more prominent.
\end{comment}

\respuesta Three changes, one per request.

\emph{JSON-LD fragment.} Figure~2 now shows a serialised component of the 4 August shadow, abridged to the fields that carry the argument: the evidence block with URL, status, bytes, latency, UTC timestamp and SHA-256; the \texttt{maltg:mapsTo} array that constitutes $\Gamma(v)$; and the \texttt{maltg:gap} field recording what the rules sought and did not find. Evidence, anchoring and gap travel together in the serialisation, which is what makes every class in $R$ traceable to an HTTP response.

\emph{Worked example.} Section~3 now carries \emph{security posture} through the full chain in prose: which root and which four subconcepts the rubric prescribes, which two signals the scan found and why, the arithmetic $\Psi=0.40\cdot 0+0.60\cdot\tfrac{2}{4}=0.30$, then $\mathrm{sds}=24.1$ and $\delta=56.3$, and finally the position that $\delta$ assigns on the remediation agenda. We chose this dimension deliberately because its outcome is none of the three extremes already shown in Figure~1, and because it makes the central point concretely: the institution plainly does something about security, yet two-thirds of what the framework prescribes leaves no machine-readable trace.

\emph{The four-versus-ten contrast.} This is now a titled paragraph of its own in Section~5, \emph{The two headline figures do not measure the same thing---and that is the finding}, which states that the scales and dimension sets do not coincide, that the comparison is a difference of frame rather than of degree, and then draws the consequence we had previously left implicit: the six dimensions missing from the institution's own assessment are, without exception, those where the ontology finds little or no machine-processable trace. The self-assessment is not merely optimistic in value; it is silent exactly where the deficit is.

We also reworked the celestial-sphere panel of Figure~3 rather than leaving it as decoration. Each constellation now carries its exact $|S_d\cap R|/|S_d|$ count and its $\Psi$, and the legend distinguishes a root that was probed and matched from one probed without match. The panel therefore no longer restates Table~3 but shows what the table cannot: that the dark constellations are questions the scan asked and answered negatively, not regions it failed to visit. Together with the revised pipeline in panel~(a), which now begins at $\Omega$ rather than at a URL, the figure conveys how the shadow is raised from the ontology. To keep the paper within a reasonable length we compressed the prose of Sections~2--6 rather than dropping the figure.

\changed{New Figure~2 (JSON-LD); worked example in Section~3; new titled paragraph in Section~5; Figure~3 redrawn with per-constellation counts and an ontology-first pipeline; Figure~4 now described in the running text.}

\section*{3.5 Discussion of limitations and policy implications}

\begin{comment}{Editor's comment}
The reflexive finding (the institution lacks the structured data the shadow requires) is insightful but underdeveloped. It should be expanded into both a methodological limitation and a concrete recommendation for public institutions.
\end{comment}

\respuesta Developed along exactly the two axes you propose, and they turn out to pull in opposite directions, which we now say.

The \emph{methodological} consequence is a measurement floor. Any instrument of this family resolves governance only to the granularity at which the observed institution has chosen to structure its public surface. Coverage is therefore a joint property of the ecosystem and of its publication practices, and $\Psi$ conflates them: a low score cannot distinguish an institution that does not govern from one that governs without publishing structured evidence of it. No refinement of $\Gamma$ removes this bound---only the institution can, by publishing. We also flag the corollary for longitudinal use, which we had missed: adopting semantic markup would raise measured coverage with no change in underlying governance, an effect that repeated measurement must control for.

The \emph{prescriptive} consequence is that the highest-leverage remediation is also the cheapest. We name three publication acts and what each unlocks: embedding \texttt{schema.org} JSON-LD, which anchors the open data and interoperability dimensions; documenting the case-consultation API that demonstrably already exists behind the single-page application, which converts an opaque service into an auditable one; and publishing a machine-readable governance manifest declaring frameworks in use, owners and review cycles, which is the minimum artefact that would let any external observer anchor the root concepts of ITIL, COBIT and security posture that are currently unreachable. None requires reforming how the institution is governed; all three change what can be known about it from outside.

\changed{\emph{A reflexive finding} in Section~6 expanded from one paragraph to three.}

\section*{4. Minor issues}

\begin{comment}{Editor's comment}
Some figures and tables are only partially described in the available text. A few sentences are slightly over-claimed (``the central result'', ``to the best of our knowledge''). Terminology consistency (digital shadow vs.\ structural digital shadow) can be tightened. The relationship with the authors' previous MALTG papers is correctly stated, but could be made even more precise regarding what is incremental and what is new.
\end{comment}

\respuesta All four addressed.

\begin{itemize}
  \item \textbf{Figure and table description.} Every float is now introduced and read in the running text. Figure~4 in particular gains a paragraph explaining what the three series are and what the shape of the plot says---the bar series converge only at the right-hand end, where open data and interoperability sit, and diverge almost completely on the left. Table~1 is introduced by the paragraph that precedes it; Table~2 and Figure~2 likewise.
  \item \textbf{Over-claims.} ``The central result'' is gone from the abstract, replaced by a statement of what the finding bounds and what it implies. ``To the best of our knowledge'' is replaced by a claim bounded to the scope of the Section~2 comparison. We also softened the assertion that nothing in the architecture is country-specific, which the transfer protocol now qualifies precisely.
  \item \textbf{Terminology.} We now define the usage once, on first mention: \emph{structural digital shadow} signals that the artefact represents components, flows and layers rather than behaviour over time; plain \emph{shadow} is used thereafter, and no distinction between the two terms is intended. We also retained the \emph{tier} (LegalShadow) versus \emph{layer} (observed ecosystem) distinction introduced in the previous round, since both words were previously doing double duty.
  \item \textbf{Relationship to MALTG.} Now stated as an explicit inventory. Reused unchanged: the governance ontology $\Omega$---its 131 classes, ten dimensions and versioned rubric. New here: the acquisition tier and evidence protocol, the anchoring $\Gamma$, the functions $\langle\Psi,\delta\rangle$, and the empirical study. We add why the distinction matters: the earlier works validated internal consistency against baseline configurations their authors had themselves declared, whereas this study changes the object of measurement rather than refining the instrument.
\end{itemize}

\hrulefill

\vspace{6pt}
\noindent\textbf{Two points we wish to flag proactively.}
The Zenodo deposit referenced in the Data Availability Statement carries a placeholder DOI in this version; the final DOI will be supplied before publication. Second, our response to \S3.1 deliberately weakens rather than strengthens several claims about generalisation. We judged that a narrow claim supported by the evidence serves the paper better than a broad one that a second jurisdiction might contradict, but we will gladly revisit this balance if you read it differently.

\vspace{8pt}
\noindent We hope the revision meets the standard for a second round, and we thank you again for a report that improved the work.

\vspace{14pt}
\noindent Yours sincerely,\\[16pt]
\textbf{Patricio M. Paccha-Angamarca}, on behalf of all authors\\
Corresponding author --- \href{mailto:patricio.paccha@epn.edu.ec}{patricio.paccha@epn.edu.ec}

\end{document}
